\color{black}

\section{Conclusão}

Este trabalho apresentou, implementou e validou a arquitetura \textbf{DIAMOND CRISTAL CLEAN} --- uma
pipeline evolutiva dual para Detecção Acústica de Anomalias em ambientes industriais, comprovadamente
viável para Edge Computing e TinyML. A metodologia baseada em 28~experimentos com quatro fases de
desenvolvimento produziu seis achados principais:

\begin{enumerate}
    \item \textbf{HHT+UKF é o componente mais crítico:} Sua remoção custa $-7$~p.p. de pAUC@0.1 e
    causa domain shift catastrófico --- resultado inédito na literatura de AAD por ablation formal.

    \item \textbf{Modelo dual é necessário:} Um único modelo não cobre adequadamente os dois regimes
    operacionais. Protocolo~A (XGBoost, pAUC~0,94) e Protocolo~B (GMM+Gamma, pAUC~0,88) são
    complementares e ambos DCASE-Ready.

    \item \textbf{pAUC@0.1 muda o ranking:} F1-Score e AUC global não são métricas adequadas para
    sistemas de alarme industrial. A adoção do pAUC@FPR~0.1 é recomendação desta pesquisa para a
    comunidade de AAD~\cite{Nishida2025}.

    \item \textbf{Deep Learning é auxiliar, não motor:} CNN, LSTM-AE e Tiny-AST são inviáveis como
    motores de inferência em Edge. O Tiny-AST tem valor exclusivo como gerador de embeddings para XAI
    em análise offline (Cloud).

    \item \textbf{Descarte documentado é contribuição:} O registro formal de 10~técnicas descartadas
    com justificativa quantitativa (pAUC, latência, memória) contribui para a comunidade ao mapear o
    espaço de soluções inviáveis --- informação raramente publicada na literatura~\cite{Khamaisi2025}.

    \item \textbf{XAI com diagnóstico físico:} A camada XAI transforma o score de anomalia em
    evidência física interpretável (hotspots temporais, faixas espectrais, mecanismo provável),
    tornando o sistema auditável para equipes de manutenção~\cite{Ali2025, Zamorano2025}.
\end{enumerate}

O sistema atende plenamente aos três critérios de Edge Computing: pAUC@0.1~$>$~0,80, latência~$<$~50~ms
e memória~$<$~4~MB. O limite operacional documentado --- Protocolo~B falha com SNR~$<-3$~dB sem rótulos
--- fornece diretriz clara para deployment industrial.

Como trabalho futuro, a submissão formal ao DCASE~2025 Task~2 consolidará o posicionamento da
arquitetura no estado da arte internacional, e a integração com firmware de microcontrolador ARM
Cortex-M validará o Cenário~3 em hardware embarcado real~\cite{Uzel2026}.

\color{black}
